\chapter{总结与展望}

\section{工作总结}

本文围绕基于时空特征融合和稀疏表示的 GNSS NLOS 识别方法展开研究，主要完成以下工作：

（1）提出 STCA 模型，设计时空双通道特征提取架构。针对 GNSS 信号数据同时包含空间分布特性和时序演化规律的特点，设计了并行的时空特征提取路径，分别捕捉卫星空间几何关系与信号质量随时间的变化模式。

（2）实现交叉注意力融合模块与稀疏表示层。通过交叉注意力机制实现时空特征的双向信息交互与自适应融合，并利用稀疏表示层对融合后的高维特征进行紧凑编码，在降低冗余的同时增强特征的判别性。

（3）在外域场景下取得 76.02\% 的识别准确率，ROC-AUC 达到 84.07\%。实验结果表明，STCA 模型能够有效区分视距信号与非视距信号，在跨场景下具有一定的泛化能力。

（4）通过完整的消融实验验证各模块的贡献。实验分别考察了交叉注意力模块及稀疏表示层对模型性能的独立影响，为架构设计的合理性提供了实证支持。时间窗口大小敏感性分析表明，窗口大小为 8 时模型性能最优。

本文的主要贡献包括：第一，设计了时空双通道特征提取架构，使模型能够同时利用卫星空间分布和信号时序演化信息；第二，引入交叉注意力机制实现时空特征的自适应融合，相比简单特征拼接提升了 6.89 个百分点的准确率；第三，采用稀疏表示层增强特征判别性，在保持模型轻量级（约 9,000 参数）的同时实现了较好的分类性能。

\section{未来工作展望}

本研究仍存在以下局限性，未来可从以下方面继续推进：

（1）数据规模与场景有限。实验数据来自 7 个测点的静态采集，总样本量约 8 万个，地域覆盖范围有限，且均为静态场景。未来可收集更多测点的数据，涵盖不同城市、不同时间段、不同天气条件，并在车载、行人等动态场景中验证模型的适用性。

（2）外域泛化性能有待提升。当前模型在跨地点场景下的准确率较训练场景有所下降，提升模型对未见环境的适应能力是实际应用面临的主要挑战。未来可探索域对抗训练、元学习、测试时自适应等域泛化技术，目标是在新地点无需重新训练即可获得较好的性能。

（3）时序模块建模能力尚未充分挖掘。当前采用的单向 LSTM 模块虽能捕捉基本的时序依赖关系，但对长程时间动态的建模能力有限。Transformer 架构、时间卷积网络等更强大的时序模型值得进一步研究。

（4）实时部署验证不足。虽然模型复杂度分析表明 STCA 属于轻量级模型，预计推理延迟在 5ms 以内，但尚未在实际接收机或移动设备上进行部署验证。未来可研究知识蒸馏、剪枝、量化等模型压缩技术，在保持性能的前提下进一步降低模型复杂度，满足边缘设备的部署需求。

（5）多星座融合具有研究价值。当前研究仅使用 GPS 单星座数据。未来可扩展到多星座融合（GPS、BDS、Galileo、GLONASS），利用更多卫星资源提升 NLOS 识别性能。同时，可融合视觉、惯性传感器等多源信息，提升模型的鲁棒性。

（6）可解释性仍有提升空间。尽管引入稀疏表示层增强了特征判别性，但深度学习模型的"黑箱"本质未变。未来可探索注意力可视化、梯度分析、特征重要性评估等可解释性方法，深入理解模型的决策机制，为模型改进提供指导。
